\documentclass[11pt]{article}

\usepackage[margin=0.9in]{geometry}
\usepackage{amsmath,amssymb,amsthm,mathtools,booktabs,enumitem,microtype}
\usepackage{xcolor}
\usepackage[colorlinks=true,linkcolor=blue!50!black,urlcolor=blue!50!black,citecolor=blue!50!black]{hyperref}

\theoremstyle{plain}
\newtheorem{theorem}{Theorem}
\newtheorem{lemma}[theorem]{Lemma}
\newtheorem{proposition}[theorem]{Proposition}
\newtheorem{corollary}[theorem]{Corollary}

\theoremstyle{definition}
\newtheorem{definition}[theorem]{Definition}

\theoremstyle{remark}
\newtheorem{remark}{Remark}

\DeclareMathOperator{\rank}{rank}
\DeclareMathOperator{\Var}{Var}
\DeclareMathOperator{\Cov}{Cov}
\DeclareMathOperator{\Bin}{Bin}

\title{\bfseries Rank-Constrained Capacity Equivalence for Gaussian Linear Associative Memory}
\author{}
\date{\vspace{-1.6em}%
\small
\begin{tabular}{@{}c@{}}
Target paper: A.~Giorlandino, S.~Goldt, and A.~Maillard,\\
\emph{Factual recall in linear associative memories: sharp asymptotics and mechanistic insights},\\
\href{https://arxiv.org/abs/2605.10795}{arXiv:2605.10795}, Conjecture~\texttt{conjecture-capacity}\\[0.35em]
Companion verification script: \texttt{solutions/linear\_assoc\_memory/explore.py}\\[0.35em]
November 2026
\end{tabular}}

\begin{document}

\maketitle

\begin{abstract}
We prove the Capacity Equivalence Conjecture (\texttt{conjecture-capacity}) of Giorlandino, Goldt, and Maillard (\href{https://arxiv.org/abs/2605.10795}{arXiv:2605.10795}) for every fixed rank fraction $0 < \kappa \le 1$: in the proportional-rank high-dimensional limit with $m/d \to \kappa$ and $p \log p / d^2 \to \alpha \in (0, \infty)$, the original shared-output associative memory problem (\textup{OP}) and the decoupled problem (\textup{DP}) share the sharp capacity threshold $\alpha_c^{\mathrm{OP}}(\kappa) = \alpha_c^{\mathrm{DP}}(\kappa) = E_\kappa / 2$, where $E_\kappa$ is the top-$\kappa$ second-moment energy of the quarter-circle law.
Since achievability below $E_\kappa / 2$ for both problems and the \textup{DP} converse were already established in \href{https://arxiv.org/abs/2605.10795}{arXiv:2605.10795}, our main result is the matching rank-constrained \textup{OP} converse for all $\alpha > E_\kappa / 2$.
Splitting any candidate $W = H + T$ into a low-rank head of rank $k = \lceil (\log p)^{24} \rceil$ and a small-operator-norm tail $T$, we prove a uniform lower bound on the shared-output competitor score $S(W)$ via Gaussian concentration, weakly correlated competitor extraction, Slepian's comparison lemma~\cite{slepian1962}, and net transfer, and contrast it with a simultaneous rank-$m$ target-energy upper bound from von Neumann's trace inequality~\cite{vonneumann1937}.
\end{abstract}

\section{Model and Main Result}

Adopt the standard-normal scaling of the main text of \href{https://arxiv.org/abs/2605.10795}{arXiv:2605.10795} (rescaling all embeddings by $d^{-1/2}$ yields the appendix scaling without changing feasibility). Let the $p$ keys $e_1, \dots, e_p \in \mathbb{R}^d$ and $p$ matched outputs $u_1, \dots, u_p \in \mathbb{R}^d$ be mutually independent $\mathcal{N}(0, I_d)$ vectors, with the injective association relabeled to the identity.
\begin{itemize}[leftmargin=2em,itemsep=2pt]
\item \textbf{Original Problem (\textup{OP}):} There exists $W \in \mathbb{R}^{d \times d}$ with $\rank(W) \le m$ such that $u_i^\top W e_i > u_j^\top W e_i$ for every $i \in [p]$ and every $j \neq i$ (ties fail).
\item \textbf{Decoupled Problem (\textup{DP}):} Each query $i \in [p]$ competes against $p - 1$ fresh $\mathcal{N}(0, I_d)$ outputs, independent across queries and independent of all keys and targets.
\end{itemize}
For $\kappa \in (0, 1]$, define $q_\kappa \in [0, 2)$ and the top-$\kappa$ quarter-circle energy $E_\kappa > 0$ by
\begin{equation}\label{eq:Ekappa}
\int_{q_\kappa}^{2} \frac{\sqrt{4 - s^2}}{\pi}\,ds = \kappa,
\qquad
E_\kappa := \int_{q_\kappa}^{2} s^2 \frac{\sqrt{4 - s^2}}{\pi}\,ds.
\end{equation}

\begin{theorem}[Capacity equivalence for all $0 < \kappa \le 1$]\label{thm:capacity}
Let $d, p, m \to \infty$ with $m/d \to \kappa \in (0, 1]$ and $p \log p / d^2 \to \alpha \in (0, \infty)$. Then
\[
\alpha_c^{\mathrm{OP}}(\kappa) = \alpha_c^{\mathrm{DP}}(\kappa) = \frac{1}{2} E_\kappa.
\]
That is, the probability that a matrix $W$ of rank at most $m$ solves \textup{OP} (resp.\ \textup{DP}) converges to $1$ if $\alpha < E_\kappa / 2$ and to $0$ if $\alpha > E_\kappa / 2$.
\end{theorem}

Giorlandino, Goldt, and Maillard~\cite{giorlandino2026} prove achievability below $E_\kappa / 2$ for both \textup{OP} and \textup{DP}, as well as impossibility above $E_\kappa / 2$ for \textup{DP}; Barnfield et~al.~\cite{barnfield2026} prove the full-rank ($\kappa = 1$) \textup{OP} converse at $\alpha > 1/2$. Proving Theorem~\ref{thm:capacity} therefore reduces to proving the \textup{OP} converse for all $\kappa \in (0, 1]$ and $\alpha > E_\kappa / 2$.

\section{Proof of the Shared-Output Rank-Constrained Converse}

Write $L := \log p$, $C := \sum_{i=1}^p u_i e_i^\top$, and
\[
S(W) := \sum_{i=1}^p \max_{j \neq i} u_j^\top W e_i.
\]
Feasibility in \textup{OP} implies $\langle C, W \rangle_F > S(W)$, and by positive homogeneity we restrict to $\|W\|_F = 1$. Throughout, $d^2 \asymp p L$, so $d \asymp \sqrt{p L}$ and every fixed power of $L$ is $o(d)$.

Set $k := \lceil L^{24} \rceil$. For any unit-Frobenius matrix $W$, split its SVD as $W = H + T$ where $H$ retains the top $k$ singular values, and set $b(W) := \|T\|_F \in [0, 1]$. Then $b$ is $1$-Lipschitz in Frobenius norm (being the Frobenius distance to the cone of matrices of rank at most $k$) and $\|T\|_{\mathrm{op}} \le 1/\sqrt{k} \le L^{-12}$.
We first establish that for every fixed $\varepsilon \in (0, 1)$, with probability tending to $1$, uniformly over all $\|W\|_F = 1$:
\begin{equation}\label{eq:unif}
\frac{S(W)}{p\sqrt{2L}} \ge c_0 := \frac{1}{4\sqrt{\pi}},
\qquad
\frac{S(W)}{p\sqrt{2L}} \ge (1 - \varepsilon)\, b(W) - o(1).
\end{equation}

\subsection*{Step 1. Two elementary properties of the output cloud}

Let $U := [u_1, \dots, u_p] \in \mathbb{R}^{d \times p}$. We claim that with probability $1 - o(1)$, every unit vector $v \in S^{d-1}$ has at least $p/3$ outputs with $v^\top u_j \ge 0$ and at least two outputs with $v^\top u_j \ge \sqrt{L}/2$.

Take a $d^{-2}$-net of $S^{d-1}$ of cardinality $\exp(O(d \log d))$. For a fixed net point $v$, the count of $j$ with $v^\top u_j \ge 0.1$ is $\Bin(p, q)$ with $q > 0.45$, so fewer than $p/3$ occurs with probability $e^{-\Omega(p)}$. At cutoff $\sqrt{L}/2 + 1$, Gaussian tail estimates give expected count $p^{7/8 + o(1)}$, so fewer than two occurs with probability at most $\exp(-p^{7/8 + o(1)})$. Since $d \log d \asymp \sqrt{p} L^{3/2}$, both failure exponents dominate $O(d \log d)$. On the high-probability event $\max_j \|u_j\| \le 2\sqrt{d}$, the net perturbation is at most $2 d^{-3/2} = o(1)$, proving both claims uniformly.

Consequently, even excluding $j = i$, $\max_{j \neq i} u_j^\top W e_i \ge \|W e_i\|\sqrt{L}/2 \ge 0$, so every summand of $S(W)$ is nonnegative. By the ``Uniform Gaussian input-frame bounds'' lemma of~\cite{giorlandino2026}, $p^{-1}\sum_{i=1}^p \|W e_i\| \ge \|W\|_F / \sqrt{2\pi}$ simultaneously for all $W$ with high probability, proving the first bound in~\eqref{eq:unif} with $c_0 = \frac{1}{2\sqrt{2}} \cdot \frac{1}{\sqrt{2\pi}} = \frac{1}{4\sqrt{\pi}}$.

\subsection*{Step 2. Weakly correlated competitors in every large subset}

\begin{lemma}\label{lem:extract}
Fix a deterministic matrix $B \in \mathbb{R}^{d \times d}$ with $\|B\|_F = 1$ and $\|B\|_{\mathrm{op}} \le L^{-10}$. Set $z_j := B^\top u_j$, $\eta := L^{-1}$, $\rho := L^{-4}$, and $n := \lfloor p/16 \rfloor$. There exists a constant $c > 0$ such that, except on an event of probability at most $\exp(-c p L^{12})$, every subset of $[p]$ of size at least $p/4$ contains $n$ indices satisfying
\begin{equation}\label{eq:compat}
1 - \eta \le \|z_j\| \le 1 + \eta,
\qquad
\biggl|\Bigl\langle \frac{z_j}{\|z_j\|}, \frac{z_l}{\|z_l\|} \Bigr\rangle\biggr| \le \rho
\qquad (j \neq l).
\end{equation}
\end{lemma}

\begin{proof}
For a fresh $z = B^\top u$ ($u \sim \mathcal{N}(0, I_d)$), we have $\mathbb{E}\|z\|^2 = \|B\|_F^2 = 1$. Since $u \mapsto \|B^\top u\|$ is $\|B\|_{\mathrm{op}}$-Lipschitz, the Gaussian Poincar\'e inequality gives $\Var(\|z\|) \le \|B\|_{\mathrm{op}}^2 \le L^{-20}$. Therefore
\[
1 - L^{-20} \le (\mathbb{E}\|z\|)^2 \le 1 \implies \mathbb{E}\|z\| \ge 1 - L^{-20}.
\]
By Gaussian Lipschitz concentration around $\mathbb{E}\|z\|$, for large $p$ the event $|\|z\| - 1| > \eta$ implies $|\|z\| - \mathbb{E}\|z\|| > \eta/2$, which has probability at most $2\exp(-c' L^{18})$. Moreover, for each fixed unit vector $v$, $\langle v, z \rangle = \langle Bv, u \rangle$ is centered Gaussian with variance $\|Bv\|^2 \le L^{-20}$. Hence the probability that a norm-good fresh vector $z$ has $|\langle v, z/\|z\| \rangle| > \rho$ for some $v$ among at most $p$ fixed unit directions is at most
\[
2p \exp\!\Bigl(-\frac{(1 - \eta)^2 \rho^2 L^{20}}{2}\Bigr).
\]
Combining both bounds, the fresh-point failure probability is at most $\exp(-c_1 L^{12})$ for large $p$.

If~\eqref{eq:compat} fails on some subset $A \subseteq [p]$ with $|A| \ge p/4$, choose an inclusion-maximal subset $J \subseteq A$ of norm-good indices satisfying the pairwise bound in~\eqref{eq:compat} (possibly $J = \varnothing$). Then $|J| < n \le p/16$, so at least $p/4 - p/16 \ge p/8$ indices in $A \setminus J$ are either norm-bad or incompatible with $J$. For each fixed pair $(J, F)$ with $F \subseteq [p] \setminus J$ and $|F| = \lceil p/8 \rceil$, condition on $\{u_l : l \in J\}$: the outputs $\{u_j : j \in F\}$ are independent, and each fails the fresh-point test relative to $\{z_l / \|z_l\| : l \in J\}$ with probability at most $\exp(-c_1 L^{12})$. A union bound over the at most $4^p$ choices of $(J, F)$ gives total failure probability at most $4^p \exp(-c_1 (p/8) L^{12}) \le \exp(-c p L^{12})$.
\end{proof}

\subsection*{Step 3. Conditional lower-tail bound via Slepian's lemma and net transfer}

Take a deterministic $d^{-3}$-net $\mathcal{N}$ of the unit Frobenius sphere in $\mathbb{R}^{d \times d}$; its log cardinality is $O(d^2 \log d) = O(p L^2)$. For each $V \in \mathcal{N}$ with $b = b(V) \ge L^{-1}$, fix an SVD split $V = H + T$ and apply Lemma~\ref{lem:extract} to $B := T/b$ (which has $\|B\|_F = 1$ and $\|B\|_{\mathrm{op}} \le L^{-11} \le L^{-10}$). Since $p L^{12} \gg p L^2$, a union bound ensures that Step~1 and Lemma~\ref{lem:extract} hold simultaneously for all such $V \in \mathcal{N}$ with probability $1 - o(1)$. Condition on a realization of $U$ in this event.

Fix one such $V \in \mathcal{N}$ and query $i \in [p]$. Split $e_i$ into independent Gaussian projections onto the top-$k$ right singular subspace of $V$ and its orthogonal complement, and condition on the first projection. This fixes $H e_i$ while leaving the orthogonal key coordinates seen by $T$ standard Gaussian and independent. By Step~1, at least $p/3 - 1 \ge p/4$ indices $j \neq i$ satisfy $u_j^\top H e_i \ge 0$ (including when $H e_i = 0$). Using only the conditioned data (e.g., lexicographically first), pick $n$ such indices $I_i$ satisfying~\eqref{eq:compat}. For $j \in I_i$, the random variables
\[
X_j := \frac{u_j^\top T e_i}{b \|z_j\|} = \Bigl\langle \frac{z_j}{\|z_j\|}, e_i \Bigr\rangle
\]
are jointly centered Gaussian with $\Var(X_j) = 1$ and $\Cov(X_j, X_l) = \langle z_j/\|z_j\|, z_l/\|z_l\| \rangle \le \rho$ ($j \neq l$), since every $z_j$ lies in the complementary right singular subspace and $I_i$ is independent of the complementary key coordinates.

Fix $\delta \in (0, 1/4)$ and set $t := (1 - \delta)\sqrt{2L}$. Let $Y_r := \sqrt{\rho}\, G_0 + \sqrt{1 - \rho}\, G_r$ ($1 \le r \le n$) with i.i.d.\ standard normals $G_0, G_1, \dots, G_n$, so $\Var(Y_r) = 1$ and $\Cov(Y_r, Y_s) = \rho \ge \Cov(X_j, X_l)$. By Slepian's Gaussian comparison lemma~\cite{slepian1962},
\[
\Pr\Bigl(\max_{j \in I_i} X_j < t \;\Big|\; U, H e_i\Bigr) \le \Pr\Bigl(\max_{1 \le r \le n} Y_r < t\Bigr).
\]
If $\sqrt{\rho}\, G_0 \ge -\frac{\delta}{4}\sqrt{2L}$, then for large $p$ the event $\max_{r \le n} Y_r < t$ implies $\max_{r \le n} G_r < (1 - \delta/2)\sqrt{2L}$. The excluded shift event has probability at most $\exp(-c_\delta L/\rho) = \exp(-c_\delta L^5)$, while the independent-maximum event has probability at most $\exp(-p^{1 - (1-\delta/2)^2 + o(1)}) \ll \exp(-c_\delta L^5)$. Because $u_j^\top H e_i \ge 0$ and $\|z_j\| \ge 1 - \eta$ on $I_i$, integrating over $H e_i$ gives, for all large $p$,
\begin{equation}\label{eq:query-tail}
\Pr\!\left(\max_{j \neq i} u_j^\top V e_i < b(1 - \eta)(1 - \delta)\sqrt{2L} \;\middle|\; U\right)
\le \exp(-c_\delta L^5).
\end{equation}

Conditional on $U$, the $p$ failure events in~\eqref{eq:query-tail} depend on distinct independent keys $e_i$ and are therefore independent. Having more than $\delta p$ failed queries has probability at most $2^p \exp(-c_\delta \delta p L^5)$, which pays for a second union bound over $|\mathcal{N}| = \exp(O(p L^2))$. By nonnegativity of the summands of $S(V)$ (Step~1), summing over the successful queries gives $S(V)/(p\sqrt{2L}) \ge (1 - \delta)^2(1 - \eta) b(V)$ simultaneously for all $V \in \mathcal{N}$ with $b(V) \ge L^{-1}$ (and with an additive $-L^{-1}$ for $b(V) < L^{-1}$).

With high probability $\max_i \|e_i\| \le 2\sqrt{d}$ as well as $\max_j \|u_j\| \le 2\sqrt{d}$, on which $|S(W) - S(V)| \le 4pd\|W - V\|_F$. Transferring from $\mathcal{N}$ (mesh $d^{-3}$) to the full unit Frobenius sphere adds $o(1)$ error after dividing by $p\sqrt{2L}$. Using $1$-Lipschitzness of $b$ and choosing $\delta$ small enough for the given $\varepsilon$ proves~\eqref{eq:unif}.

\subsection*{Step 4. Comparing target energy to competitor scores}

In standard-normal units, the correlation-matrix estimates of~\cite{giorlandino2026} give
\[
\frac{\|C\|_{\mathrm{op}}}{\sqrt{pd}} = O_{\Pr}(1),
\qquad
\frac{\|C_m\|_F}{d\sqrt{p}} \xrightarrow{\Pr} \sqrt{E_\kappa},
\]
where $C_m$ is the rank-$m$ truncated SVD of $C$. (Conditional on $U$, the rescaled matrix $C/\sqrt{pd}$ has the law of $(U U^\top/p)^{1/2} G/\sqrt{d}$ with $G$ i.i.d.\ standard Gaussian; since $p/d \to \infty$, we have $U U^\top / p \to I_d$ in operator norm, and the quarter-circle law yields the limit $\sqrt{E_\kappa}$.) Crucially, since $\kappa \in (0, 1]$, we have $q_\kappa < 2$ and therefore $E_\kappa > 0$.

For every unit-Frobenius matrix $W = H + T$ with $\rank(W) \le m$, we have $\|H\|_* \le \sqrt{k}\|H\|_F \le \sqrt{k}$, $\rank(T) \le m$, and $\|T\|_F = b(W)$. By von Neumann's trace inequality~\cite{vonneumann1937}, simultaneously for all such $W$:
\begin{equation}\label{eq:target-ub}
\frac{\langle C, W \rangle_F}{d\sqrt{p}}
\le \frac{\sqrt{k}\|C\|_{\mathrm{op}}}{d\sqrt{p}} + \frac{\|C_m\|_F}{d\sqrt{p}}\, b(W)
\le \sqrt{E_\kappa}\, b(W) + o_{\Pr}(1),
\end{equation}
because $\sqrt{k/d} = O(L^{12}/\sqrt{d}) \to 0$. No independence between~\eqref{eq:unif} and~\eqref{eq:target-ub} is needed.

Now assume $\alpha > E_\kappa / 2$, and pick $\varepsilon \in (0, 1)$ with $(1 - \varepsilon)\sqrt{2\alpha} > \sqrt{E_\kappa}$. Converting the denominator in~\eqref{eq:unif} from $p\sqrt{2L}$ to $d\sqrt{p}$ multiplies by $\sqrt{2pL}/d \to \sqrt{2\alpha}$. If a feasible unit-Frobenius matrix $W$ of rank at most $m$ exists, then $\langle C, W \rangle_F > S(W)$; combining the coarse bound in~\eqref{eq:unif} with~\eqref{eq:target-ub} yields
\[
c_0 \sqrt{2\alpha} - o_{\Pr}(1) \le \frac{S(W)}{d\sqrt{p}} < \frac{\langle C, W \rangle_F}{d\sqrt{p}} \le \sqrt{E_\kappa}\, b(W) + o_{\Pr}(1).
\]
Since $E_\kappa > 0$, this forces $b(W) \ge b_0 := \frac{c_0 \sqrt{2\alpha}}{2\sqrt{E_\kappa}} > 0$ with probability tending to $1$. Inserting $b(W) \ge b_0$ into the sharp bound in~\eqref{eq:unif} and comparing with~\eqref{eq:target-ub} gives
\[
0 < \frac{\langle C, W \rangle_F - S(W)}{d\sqrt{p}}
\le -\bigl((1 - \varepsilon)\sqrt{2\alpha} - \sqrt{E_\kappa}\bigr) b_0 + o_{\Pr}(1),
\]
which is strictly negative with probability tending to $1$, a contradiction. Because~\eqref{eq:unif} and~\eqref{eq:target-ub} are simultaneous in $W$ (covering data-dependent feasible solutions), $\Pr(\text{\textup{OP} feasible}) \to 0$. This proves Theorem~\ref{thm:capacity}. \hfill$\square$

\begin{remark}[Boundary $\alpha = E_\kappa / 2$ and proof structure]
The conjecture makes strict-side assertions. At the limiting ratio $\alpha = E_\kappa / 2$, diagonal sequences approaching $E_\kappa / 2$ from below (resp.\ above) yield limiting success probability $1$ (resp.\ $0$); finer critical-window asymptotics require additional rate assumptions. Notice also that standard i.i.d.\ Gordon min--max or Cover-counting arguments do not apply directly to the dependent non-Gaussian rows $e_i \otimes (u_i - u_j)$, whereas conditioning on $U$ and the top-$k$ head projection produces genuinely Gaussian vectors with controlled covariance for Slepian's lemma while rank dependence enters cleanly via $\|C_m\|_F$.
\end{remark}

\section{Finite-Dimensional Verification}

Running the script \texttt{solutions/linear\_assoc\_memory/explore.py} under \texttt{.venv/bin/python -B} (seed \texttt{260510795}) tests $d \in \{12, 16\}$ across four loads with four paired \textup{OP}/\textup{DP} trials per load using L-BFGS factor search (verifying normalized margin $> 10^{-8}$ on all $p(p-1)$ pairs) and HiGHS constraint-generation LP on full-rank failures. Table~\ref{tab:sweep} summarizes the 64 model instances ($43$ full-rank witnesses, $21$ LP-infeasible; limiting thresholds $E_\kappa / 2$ at $\kappa \in \{0.25, 0.50, 0.75, 1.00\}$ are $0.31145628$, $0.44686540$, $0.49352390$, and $0.50000000$).

\begin{table}[ht]
\centering
\caption{Checked witness counts (in rank-fraction order $\kappa = 0.25, 0.50, 0.75, 1.00$) and full-rank LP infeasibility counts out of $4$ trials per cell.}
\label{tab:sweep}
\small
\begin{tabular}{@{}rrrccc@{}}
\toprule
$d$ & $p$ & $p\log(p)/d^2$ & \textup{OP} witnesses & \textup{DP} witnesses & Full-rank LP infeasible (\textup{OP}/\textup{DP}) \\
\midrule
12 & 16 & 0.3081 & 4, 4, 4, 4 & 4, 4, 4, 4 & 0 / 0 \\
12 & 26 & 0.5883 & 4, 4, 4, 4 & 3, 4, 4, 4 & 0 / 0 \\
12 & 39 & 0.9922 & 0, 2, 4, 4 & 0, 1, 3, 3 & 0 / 1 \\
12 & 51 & 1.3925 & 0, 0, 0, 0 & 0, 0, 0, 0 & 4 / 4 \\
16 & 24 & 0.2979 & 4, 4, 4, 4 & 4, 4, 4, 4 & 0 / 0 \\
16 & 41 & 0.5948 & 4, 4, 4, 4 & 3, 4, 4, 4 & 0 / 0 \\
16 & 62 & 0.9995 & 0, 0, 2, 2 & 0, 0, 1, 2 & 2 / 2 \\
16 & 81 & 1.3904 & 0, 0, 0, 0 & 0, 0, 0, 0 & 4 / 4 \\
\bottomrule
\end{tabular}
\end{table}

\begin{thebibliography}{4}
\bibitem{slepian1962}
D.~Slepian, \emph{The One-Sided Barrier Problem for Gaussian Noise}, Bell System Technical Journal \textbf{41}(2):463--501, 1962.
\bibitem{vonneumann1937}
J.~von Neumann, \emph{Some matrix-inequalities and metrization of matric-space}, Tomsk Univ.\ Rev.\ \textbf{1}:286--300, 1937.
\bibitem{giorlandino2026}
A.~Giorlandino, S.~Goldt, and A.~Maillard, \emph{Factual recall in linear associative memories: sharp asymptotics and mechanistic insights}, \href{https://arxiv.org/abs/2605.10795}{arXiv:2605.10795}, 2026.
\bibitem{barnfield2026}
M.~Barnfield et~al., \emph{Sharp capacity thresholds for linear associative memory}, \href{https://arxiv.org/abs/2605.05189}{arXiv:2605.05189}, 2026.
\end{thebibliography}

\end{document}
